\section{Complete proof of the six-input theorem}\label{app:six}

This appendix gives the full parity-by-parity construction for
\Cref{thm:six720}.  In every product below, the displayed factors are built
from disjoint leaves.

\subsection{Six even inputs}\label{app:six0}

First construct an $18$-block.  If one input is divisible by $9$, multiply it
by another even input.  Otherwise six residues occupy the four nonzero
signed classes $\{\pm1\},\{\pm2\},\{\pm3\},\{\pm4\}$ modulo $9$, so two
collide; their suitable sum or difference is divisible by $9$ and by $2$.

Among the remaining four even inputs, construct a $10$-block.  If one is
divisible by $5$, multiply it by another even input.  Otherwise a collision
in the two nonzero signed classes modulo $5$ yields a sum or difference
divisible by $5$ and by $2$.  The final two even inputs have product divisible
by $4$.  Hence $18\cdot10\cdot4=720$.

\subsection{Six odd inputs}\label{app:six6}

Construct an $18$-block as follows.  If at least two inputs are divisible by
$9$, their sum or difference is divisible by $18$.  If exactly one is
divisible by $9$, the other five occupy four nonzero signed classes modulo
$9$; if none is divisible by $9$, all six occupy those four classes.  In both
cases a collision gives a sum or difference divisible by $9$ and by $2$.

Among the remaining four odd inputs, construct a $10$-block.  If at least two
are divisible by $5$, take their sum or difference.  If exactly one is
divisible by $5$, the other three occupy two nonzero signed classes; if none
is divisible by $5$, all four occupy those two classes.  Thus a collision
gives a sum or difference divisible by $10$.  The final odd pair yields $4$.
Therefore $18\cdot10\cdot4=720$.

\subsection{One odd and five even inputs}\label{app:six1}

If one of the five even inputs is divisible by $5$, it supplies $10$; reserve
another even input as a factor $2$, and construct a $36$-multiple from the
remaining four leaves.  This gives $10\cdot2\cdot36=720$.

Otherwise the five even inputs lie in the four nonzero even signed classes modulo $20$,\[\{\pm2\},\quad\{\pm4\},\quad\{\pm6\},\quad\{\pm8\}.\]  A collision gives a
$20$-block, while the other four leaves give $36$.  Thus
$20\cdot36=720$.

\subsection{Five odd and one even input}\label{app:six5}

Write the odd inputs as $a,b,c,d,f$ and the even input as $e$.  If
$5\mid a$, among the other four odd inputs choose $b,c$ with
$12\mid b\pm c$.  Then $60\mid a(b\pm c)$, while $d,f,e$ yield a
$12$-multiple.  Thus $60\cdot12=720$.

If no odd input is divisible by $5$, the five odd inputs occupy four nonzero
odd signed classes modulo $20$, so a pair yields $20$.  The remaining three
odd inputs yield $18$, and multiplication by the unused even input turns that
block into a multiple of $36$.  Thus $20\cdot36=720$.

\subsection{Two odd and four even inputs}\label{app:six2}

Let the odd inputs be $a,b$ and the even inputs
$e_1,e_2,e_3,e_4$.

\paragraph{A. A multiple of $5$ occurs.}
If an odd input $a$ is divisible by $5$, choose three even inputs and
construct a $48$-multiple.  The remaining odd-even pair yields a multiple of
$3$.  Hence $5\cdot48\cdot3=720$.

If no odd input is divisible by $5$ but an even input is, that even input
supplies $10$.  From the remaining five leaves, a $2\OO1\EE$ triple yields
$12$ and the final even pair yields $6$.  Thus $10\cdot12\cdot6=720$.

Assume now that all six inputs are nonzero modulo $5$.

\paragraph{B. The two odd inputs occupy one signed class modulo $5$.}
Their suitable sum or difference is divisible by $10$, while the four even
inputs yield $72$.  Thus $10\cdot72=720$.

Henceforth the two odd inputs occupy different signed classes modulo $5$.

\paragraph{C. Two even inputs occupy one signed class modulo $20$.}
Their suitable sum or difference is divisible by $20$, and the other four
leaves yield $36$.  Thus $20\cdot36=720$.

We may therefore assume that the four even inputs occupy the four distinct
signed classes
\[
 \{\pm2\},\ \{\pm4\},\ \{\pm6\},\ \{\pm8\}\pmod{20}.
\]
Group $\{\pm2\},\{\pm8\}$ as $X$ and
$\{\pm4\},\{\pm6\}$ as $Y$.  These are the two nonzero signed classes
modulo $5$, and two even inputs lie in each group.

\paragraph{D. A multiple of $3$ occurs.}
If an odd input $a$ is divisible by $3$, the other odd input $b$ has the same
modulo-$5$ signed class as two even inputs.  Choose one of them, $e$, with a
sign satisfying $5\mid b\pm e$.  The remaining three even inputs yield $48$,
so $3\cdot5\cdot48=720$.

If neither odd input is divisible by $3$ but an even input $e$ is, then $e$
supplies $6$.  The two even inputs in the modulo-$5$ group opposite to $e$
yield $10$, and the remaining two odd inputs and one even input yield $12$.
Thus $6\cdot10\cdot12=720$.

Assume finally that all six inputs are coprime to $15$.

\paragraph{E. No input is divisible by $3$ or $5$.}
The four even inputs occupy the three nonzero signed classes
$\{\pm1\},\{\pm2\},\{\pm4\}$ modulo $9$, so two yield an $18$-block.
If this pair lies wholly in $X$ or wholly in $Y$, the remaining two even
inputs lie in the other modulo-$5$ class and yield $10$, while the odd pair
yields $4$.  This gives $18\cdot10\cdot4=720$.

If the $18$-pair contains one input from each of $X,Y$, then the two unused
even inputs also contain one from each group.  Together with the two odd
inputs, which occupy the two different modulo-$5$ classes, they form a
$2\OO2\EE$ block and yield $40$.  Hence $18\cdot40=720$.

\subsection{Three odd and three even inputs}\label{app:six3}

Let $a,b,c$ be odd and $x,y,z$ even.  If one even input is divisible by $3$,
it supplies $6$, and the remaining five inputs yield $120$.  Otherwise the
three even inputs are nonzero modulo $3$ and occupy only the two residue
classes $0,2\pmod4$.  Two, say $x,y$, are congruent modulo $4$.  Choose a
sign with $3\mid x\pm y$; the same sum or difference is divisible by $4$,
hence by $12$.  The remaining three odd inputs and one even input yield $60$.
Thus $12\cdot60=720$.

\subsection{Four odd and two even inputs}\label{app:six4}

Let $a,b,c,d$ be odd and $x,y$ even.  If one even input is divisible by $3$,
it supplies $6$ and the remaining five inputs yield $120$.

Assume $3\nmid xy$.  If $x\equiv y\pmod4$, choose the sign for which
$3\mid x\pm y$; the same block is divisible by $4$, hence by $12$.  The
four odd inputs yield $60$, giving $12\cdot60=720$.

It remains that the two even inputs have different residues modulo $4$.
Relabel so that $4\mid x$ and $y\equiv2\pmod4$.
If an odd input $a$ is divisible by $3$, then $12\mid ax$, and the remaining
three odd inputs and one even input yield $60$.

Suppose no odd input is divisible by $3$, but one odd input $a$ is divisible
by $5$.  Then $20\mid ax$, while the remaining four inputs yield $36$.

Finally assume all four odd inputs are coprime to $15$.  They yield $120$ by
\Cref{thm:4ocop120}.  Since $x,y$ are even and nonzero modulo $3$, their
suitable sum or difference is divisible by $6$.  Thus $120\cdot6=720$.

The seven subsections exhaust all parity patterns.
